\let\negmedspace\undefined
\let\negthickspace\undefined
\documentclass[journal]{IEEEtran}
\usepackage[a5paper, margin=10mm, onecolumn]{geometry}
%\usepackage{lmodern} % Ensure lmodern is loaded for pdflatex
\usepackage{tfrupee} % Include tfrupee package

\setlength{\headheight}{1cm} % Set the height of the header box
\setlength{\headsep}{0mm}     % Set the distance between the header box and the top of the text

\usepackage{gvv-book}
\usepackage{gvv}
\usepackage{cite}
\usepackage{amsmath,amssymb,amsfonts,amsthm}
\usepackage{algorithmic}
\usepackage{graphicx}
\usepackage{textcomp}
\usepackage{xcolor}
\usepackage{txfonts}
\usepackage{listings}
\usepackage{enumitem}
\usepackage{mathtools}
\usepackage{gensymb}
\usepackage{comment}
\usepackage[breaklinks=true]{hyperref}
\usepackage{tkz-euclide} 
\usepackage{listings}
% \usepackage{gvv}                                        
\def\inputGnumericTable{}                                 
\usepackage[latin1]{inputenc}                                
\usepackage{color}                                            
\usepackage{array}                                            
\usepackage{longtable}                                       
\usepackage{calc}                                             
\usepackage{multirow}                                         
\usepackage{hhline}                                           
\usepackage{ifthen}                                           
\usepackage{lscape}

\usepackage{circuitikz}

\raggedbottom
\begin{document}

\bibliographystyle{IEEEtran}
\vspace{3cm}

\title{2015-EE-27-39}
\author{EE24BTECH11065-spoorthi
}
% \maketitle
% \newpage
% \bigskip
{\let\newpage\relax\maketitle}

\renewcommand{\thefigure}{\theenumi}
\renewcommand{\thetable}{\theenumi}
\setlength{\intextsep}{10pt} % Space between text and floats


\numberwithin{equation}{enumi}
\numberwithin{figure}{enumi}
\renewcommand{\thetable}{\theenumi}

\begin{enumerate}

\item The closed loop line integral
\[
\oint_{|z|=5} \frac{z^3 + z^2 + 8}{z + 2} \, dz
\]
evaluated counter-clockwise, is \hfill{[EE-2019]}
\begin{multicols}{4}
    

\begin{enumerate}
    \item  $+8j\pi$
    \item  $-8j\pi$
    \item  $-4j\pi$
    \item  $+4j\pi$
\end{enumerate}
\end{multicols}

\item A periodic function $f(t)$, with a period of $2\pi$, is represented as its Fourier series,
\[
f(t) = a_0 + \sum_{n=1}^{\infty} a_n \cos nt + \sum_{n=1}^{\infty} b_n \sin nt.
\]
If
\[
f(t) = \begin{cases} 
      A \sin t, & 0 \leq t \leq \pi \\
      0, & \pi < t < 2\pi 
   \end{cases}
\]
the Fourier series coefficients $a_1$ and $b_1$ of $f(t)$ are\hfill{[EE-2019]}
\begin{multicols}{2}
    

\begin{enumerate}
    \item  $a_1 = \frac{A}{4}; \, b_1 = 0$
    \item  $a_1 = \frac{A}{4}; \, b_1 = 0$
    \item  $a_1 = 0; \, b_1 = \frac{A}{\pi}$
    \item  $a_1 = 0; \, b_1 = \frac{4}{2}$
\end{enumerate}
\end{multicols}
\item The asymptotic Bode magnitude plot of a minimum phase transfer function $G(s)$ is shown below.
\begin{figure}[H]
\centering

\resizebox{0.5\textwidth}{!}{%
\begin{circuitikz}
\tikzstyle{every node}=[font=\Large]
\draw [line width=0.8pt, ->, >=Stealth] (10.75,17.75) -- (10.75,25.5);
\draw [line width=0.8pt, ->, >=Stealth] (13.75,18.25) -- (25.75,18.25);
\draw [line width=0.8pt, short] (10.5,23.5) -- (11,23.5);
\draw [line width=0.8pt, short] (10.5,21.75) -- (11,21.75);
\draw [line width=0.8pt, short] (10.5,20) -- (11,20);
\draw [line width=0.8pt, short] (10,18.25) -- (14.75,18.25);
\draw [line width=0.8pt, short] (14.75,18) -- (14.75,18.5);
\draw [line width=0.8pt, short] (18.5,18.5) .. controls (18.5,18) and (18.5,18.25) .. (18.5,18);
\draw [line width=0.8pt, short] (22,18.5) -- (22,18);
\draw [line width=0.8pt, dashed] (10.75,21.75) -- (14.5,21.75);
\draw [line width=0.8pt, dashed] (14.75,21.75) -- (14.75,18.25);
\draw [line width=1.4pt, short] (14.75,21.75) -- (11,23);
\draw [line width=1.4pt, short] (14.75,21.75) -- (18.5,18.25);
\draw [line width=1.4pt, short] (18.5,18.25) -- (22,16.25);
\draw [line width=1.4pt, short] (22,16.25) -- (23,14.5);
\draw [line width=0.7pt, dashed] (22,18) -- (22,16.25);
\node [font=\normalsize] at (14.75,17.75) {1};
\node [font=\normalsize] at (18.5,17.75) {10};
\node [font=\normalsize] at (22.5,17.75) {20};
\node [font=\normalsize] at (23.75,17.75) {$\omega$};
\node [font=\normalsize] at (10.25,18.5) {0};
\node [font=\normalsize] at (10,20) {20};
\node [font=\normalsize] at (10,21.75) {40};
\node [font=\normalsize] at (10,23.5) {60};

\node [font=\normalsize] at (12.25,24.5) {$\mid G(j\omega)\mid$ (dB)};
\node [font=\normalsize] at (24.75,17.75) {(log scale)};
\node [font=\normalsize] at (14.5,22.75) {-20 dB/decade};
\node [font=\normalsize] at (18.75,20.25) {-40 dB/decade};
\node [font=\normalsize] at (24,15.5) {-60 dB/decade};
\end{circuitikz}
}%




\end{figure}

Consider the following two statements.
\begin{itemize}
    \item Statement I: Transfer function $G(s)$ has three poles and one zero.
    \item Statement II: At very high frequency ($\omega \to \infty$), the phase angle $\angle G(j\omega) = -\frac{3\pi}{2}$.
\end{itemize}

Which one of the following options is correct?\hfill{[EE-2019]}
\begin{enumerate}
    \item  Statement I is true and statement II is false.
    \item  Statement I is false and statement II is true.
    \item  Both the statements are true.
    \item  Both the statements are false.
\end{enumerate}

\item The transfer function of a phase lead compensator is given by
\[
D(s) = \frac{3(s+\frac{1}{3T})}{s+\frac{1}{T}}.
\]
The frequency (in rad/sec), at which $\angle D(j\omega)$ is maximum, is\hfill{[EE-2019]}
\begin{multicols}{4}
    

\begin{enumerate}
    \item  $\sqrt{\frac{3}{T^2}}$
    \item  $\sqrt{\frac{1}{3T^2}}$
    \item  $\sqrt{3T}$
    \item  $\sqrt{3T^2}$
\end{enumerate}
\end{multicols}
\item Consider a state-variable model of a system
\[
\dot{x}_1 = \begin{bmatrix} 0 & 1 \\ -\alpha & -2\beta \end{bmatrix} x_2 + \begin{bmatrix} 0 \\ \alpha \end{bmatrix} r
\]
where $y$ is the output, and $r$ is the input. The damping ratio $\xi$ and the undamped natural frequency $\omega_n$ $\brak{rad/sec}$ of the system are given by \hfill{[EE-2019]}
\begin{multicols}{2}
    

\begin{enumerate}
    \item  $\xi = \frac{\beta}{\sqrt{\alpha}}; \, \omega_n = \sqrt{\alpha}$
    \item  $\xi = \sqrt{\alpha}; \, \omega_n = \frac{\beta}{\sqrt{\alpha}}$
    \item  $\xi = \frac{\beta}{\alpha}; \, \omega_n = \sqrt{\beta}$
    \item  $\xi = \sqrt{\beta}; \, \omega_n = \sqrt{\alpha}$
\end{enumerate}
\end{multicols}
\item A moving coil instrument having a resistance of $10 \, \Omega$, gives a full-scale deflection when the current is $10 \, \text{mA}$. What should be the value of the series resistance, so that it can be used as a voltmeter for measuring potential difference up to $100 \, \text{V}$? \\. \hfill{[EE-2019]}
\begin{multicols}{4}
    

\begin{enumerate}
    \item  $9 \, \Omega$
    \item  $99 \, \Omega$
    \item  $990 \, \Omega$
    \item  $9990 \, \Omega$
\end{enumerate}
\end{multicols}


\item The enhancement type MOSFET in the circuit below operates according to the square law. $\mu_n C_{\text{ox}} = 100 \, \mu \text{A}/\text{V}^2$, the threshold voltage ($V_T$) is 500 mV. Ignore channel length modulation. The output voltage $V_{\text{out}}$ is \hfill{[EE-2019]}
\begin{figure}[H]
\centering


\resizebox{0.3\textwidth}{!}{%
\begin{circuitikz}
\tikzstyle{every node}=[font=\footnotesize]
\draw [ line width=0.5pt](11.75,22) to[short] (13.5,22);
\draw [ line width=0.5pt](12.75,22) to[short] (12.75,19.5);
\draw [ line width=0.5pt](12.75,19.5) to[american current source] (12.75,16.75);
\draw [ line width=0.5pt](10.75,16.75) to[short, -o] (14.75,16.75) ;
\node at (12.75,16.75) [circ] {};
\draw [ line width=0.5pt](12.75,16.75) to[short] (12.75,15.75);
\draw [ line width=0.5pt](10.75,16.75) to[short] (10.75,15);
\draw [ line width=0.5pt](12.75,15.75) to[short] (11.75,15.75);
\draw [ line width=0.5pt](11.75,15.75) to[short] (11.75,14.25);
\draw [ line width=0.5pt](11.75,14.25) to[short] (12.75,14.25);
\draw [ line width=0.5pt](12.75,14.25) to[short] (12.75,13.5);
\draw [ line width=0.5pt](10.75,15) to[short] (11.5,15);
\draw [ line width=0.5pt](11.5,15.75) to[short] (11.5,14.25);
\draw [line width=0.5pt, ->, >=Stealth] (12.25,14.25) -- (12.5,14.25);
\draw [line width=0.5pt](12.75,13.5) to (12.75,12.5) node[ground]{};
\node [font=\small] at (12.5,22.5) {V};
\node [font=\footnotesize] at (12.75,22.25) {DO};
\node [font=\small] at (13.75,18) {5 $\mu A$};
\node [font=\small] at (13.5,15.5) {W};
\node [font=\small] at (13.5,15) {L};
\node [font=\small] at (14.25,15.25) {=};
\node [font=\small] at (15.25,15.5) {$10\mu m$};
\node [font=\small] at (15.25,15) {$1\mu m$};
\draw [line width=0.5pt, short] (13.25,15.25) -- (13.75,15.25);
\draw [line width=0.5pt, short] (14.75,15.25) -- (16,15.25);
\node [font=\normalsize] at (15.5,16.75) {V};
\node [font=\footnotesize] at (15.75,16.5) {out};
\end{circuitikz}
}%



\end{figure}
\begin{multicols}{4}
    


\begin{enumerate}
    \item  100 $mV$
    \item  500 $mV$
    \item  600 $mV$
    \item  2 $V$
\end{enumerate}
\end{multicols}
\item In the circuit below, the operational amplifier is ideal. If $V_1 = 10 \, \text{mV}$ and $V_2 = 50 \, \text{mV}$, the output voltage ($V_{\text{out}}$) is \hfill{[EE-2019]}
\begin{figure}[H]
\centering

\resizebox{0.5\textwidth}{!}{%
\begin{circuitikz}
\tikzstyle{every node}=[font=\large]













\draw [short] (8,12.5) -- (8,10.75);
\draw [short] (8,12.5) -- (9.5,11.5);
\draw [short] (8,10.75) -- (9.5,11.5);
\draw [short] (9.5,11.5) -- (10,11.5);
\node [font=\LARGE] at (8.25,11.25) {\textbf{+}};
\node [font=\LARGE] at (8.25,12) {\textbf{-}};
\draw (10,11.5) to[short, -o] (12,11.5) ;
\node [font=\normalsize] at (6,12) {o};
\node [font=\normalsize] at (6,11.25) {o};
\draw (6,12) to[european resistor] (8,12);
\draw (6,11.25) to[european resistor] (8,11.25);
\draw (7.75,11.25) to[european resistor] (7.75,9);
\draw (7.75,9) to[short] (7.75,8.5);
\draw (5.75,8.75) to[short] (10,8.75);
\draw (7.75,8.5) to (7.75,8.25) node[ground]{};
\draw (7.75,13.75) to[european resistor] (10,13.75);
\draw (10,13.75) to[short] (10,11.5);
\draw (8,13.75) to[short] (7.75,13.75);
\draw (7.75,13.75) to[short] (7.75,12);
\node [font=\large] at (5.75,12) {V};
\node [font=\footnotesize] at (6,11.75) {1};
\node [font=\large] at (5.75,10.75) {V};
\node [font=\footnotesize] at (6,10.5) {2};
\node [font=\footnotesize] at (12,11.75) {out};
\node [font=\large] at (11.75,12) {V};
\node [font=\large] at (6.75,10.5) {10 k$\Omega$};
\node [font=\large] at (6.5,12.75) {10 k$\Omega$};
\node [font=\large] at (8.75,14.25) {100 k$\Omega$};
\node [font=\large] at (8.75,10) {100 k$\Omega$};
\end{circuitikz}
}%

\label{fig:my_label}

\end{figure}
\begin{multicols}{4}
    


\begin{enumerate}
    \item  100 $mV$
    \item  400 $mV$
    \item  500 $mV$
    \item  600 $mV$
\end{enumerate}
\end{multicols}



    \item The output expression for the Karnaugh map shown below is \hfill{[EE-2019]}
    \begin{figure}[H]
\centering

\resizebox{0.5\textwidth}{!}{%
\begin{circuitikz}
\tikzstyle{every node}=[font=\Huge]
\draw (9.25,14) to[short] (9.5,14);
\draw (9.25,13.5) to[short] (9.5,13.5);
\draw (9.5,14) node[ieeestd and port, anchor=in 1, scale=0.89](port){} (port.out) to[short] (11.25,13.75);
\draw (9.25,11) to[short] (9.5,11);
\draw (9.25,10.5) to[short] (9.5,10.5);
\draw (9.5,11) node[ieeestd and port, anchor=in 1, scale=0.89](port){} (port.out) to[short] (11.25,10.75);
\draw (5.25,9.25) node[ieeestd not port, anchor=in](port){} (port.out) to[short] (7,9.25);
\draw (port.in) to[short] (5,9.25);
\draw (13.25,12.75) to[short] (13.5,12.75);
\draw (13.25,12.25) to[short] (13.5,12.25);
\draw (13.5,12.75) node[ieeestd or port, anchor=in 1, scale=0.89](port){} (port.out) to[short] (15.25,12.5);
\draw (11.25,13.75) to[short] (12.5,13.75);
\draw (12.5,13.75) to[short] (13.25,13.75);
\draw (13.25,13.75) to[short] (13.25,12.75);
\draw (11.25,10.75) to[short] (13.25,10.75);
\draw (13.25,12.25) to[short] (13.25,10.75);
\draw (7,9.25) to[short] (9.25,9.25);
\draw (9.25,10.5) to[short] (9.25,9.25);
\draw (6.75,14) to[short] (8.75,14);
\draw (6.75,14) to[short] (9.25,14);
\draw (4.75,14) to[short] (4,14);
\draw (9.25,13.5) to[short] (6.25,13.5);
\draw (6.25,13.5) to[short] (6.25,12.75);
\draw (6.25,12.75) to[short] (4,12.75);
\draw (9.25,11) to[short] (4,11);
\draw (5,14) node[ieeestd not port, anchor=in](port){} (port.out) to[short] (6.75,14);
\draw (port.in) to[short] (4.75,14);
\draw (4.5,14) to[crossing] (4.5,11.25);
\draw (4.75,12.75) to[crossing] (4.75,9.25);
\draw (5,9.25) to[short] (4.5,9.25);
\node at (4.75,12.75) [circ] {};
\node at (4.5,14) [circ] {};
\node [font=\Large] at (4.25,14.25) {X};
\node [font=\Large] at (4.25,13) {\textbf{Y}};
\node [font=\Large] at (15.25,12.75) {Z};
\draw (15.25,12.5) to[short] (16,12.5);
\draw (4.5,11.25) to[short] (4.5,11);
\draw [short] (4.5,11.25) -- (4.5,11);
\end{circuitikz}
}%
\label{fig:my_label}

\end{figure}
\begin{multicols}{2}
    

    \begin{enumerate}
        \item $\overline{Q}R + S$
        \item $Q\overline{R} + \overline{S}$
        \item $QR + S$
        \item $QR + \overline{S}$
    \end{enumerate}
    \end{multicols}
    \item In the circuit shown below, $X$ and $Y$ are digital inputs, and $Z$ is a digital output. The equivalent circuit is a \hfill{[EE-2019]}
    \begin{figure}[H]
\centering
    
\resizebox{0.5\textwidth}{!}{%
\begin{circuitikz}
\tikzstyle{every node}=[font=\Large]
\draw [line width=0.8pt, short] (11.5,22) -- (16.5,22);
\draw [line width=0.8pt, short] (16.5,22) -- (16.5,17);
\draw [line width=0.8pt, short] (11.5,17) -- (16.5,17);
\draw [line width=0.8pt, short] (11.5,22) -- (11.5,17);
\draw [line width=0.8pt, short] (12.75,22) -- (12.75,17);
\draw [line width=0.8pt, short] (14,22) -- (14,17);
\draw [line width=0.8pt, short] (15.25,22) -- (15.25,17);
\draw [line width=0.8pt, short] (11.5,20.75) -- (16.5,20.75);
\draw [line width=0.8pt, short] (11.5,19.5) -- (16.5,19.5);
\draw [line width=0.8pt, short] (11.5,18.25) -- (16.5,18.25);
\draw [line width=0.8pt, short] (11.5,22) -- (10,23.5);
\node [font=\Large] at (12,21.25) {0};
\node [font=\Large] at (12,17.5) {0};
\node [font=\Large] at (12,20.25) {};
\node [font=\Large] at (14.5,17.5) {0};
\node [font=\Large] at (16,17.75) {};
\node [font=\Large] at (16,21.25) {0};
\node [font=\Large] at (13.25,21.25) {1};
\node [font=\Large] at (14.5,21.25) {1};
\node [font=\Large] at (12.25,20) {1};
\node [font=\Large] at (13.25,20) {1};
\node [font=\Large] at (14.5,20) {1};
\node [font=\Large] at (16,20) {1};
\node [font=\Large] at (12,18.75) {1};
\node [font=\Large] at (13.25,18.75) {1};
\node [font=\Large] at (14.5,18.75) {1};
\node [font=\Large] at (16,18.75) {1};
\node [font=\Large] at (12,22.5) {00};
\node [font=\Large] at (11,21.25) {00};
\node [font=\Large] at (13.5,22.5) {01};
\node [font=\Large] at (11,20) {01};
\node [font=\Large] at (14.75,22.5) {11};
\node [font=\Large] at (11,18.75) {11};
\node [font=\Large] at (16,22.5) {10};
\node [font=\Large] at (11,17.5) {10};
\node [font=\Large] at (11.25,23) {PQ};
\node [font=\Large] at (10.25,22.25) {RS};
\node[font=\Large]  at (16,17.5) {0};
\node[font=\Large]  at (13.4,17.5) {0};
\draw [line width=0.8pt, short] (16.5,19.75) -- (16.5,19.5);
\end{circuitikz}
}%

\label{fig:my_label}

    \end{figure}
    \begin{multicols}{4}
        
    
    \begin{enumerate}
        \item NAND gate
        \item NOR gate
        \item XOR gate
        \item XNOR gate
    \end{enumerate}
    \end{multicols}
    \item A DC-DC buck converter operates in continuous conduction mode. It has 48 $V$ input voltage, and it feeds a resistive load of 24 $\Omega$. The switching frequency of the converter is 250 $Hz$. If switch-on duration is 1 $ms$, the load power is \hfill{[EE-2019]}
    \begin{multicols}{4}
        
    
    \begin{enumerate}
        \item 6 $W$
        \item 12 $W$
        \item 24 $W$
        \item 48 $W$
    \end{enumerate}
    \end{multicols}
    \item The line currents of a three-phase four wire system are square waves with amplitude of 100 $A$. These three currents are phase shifted by 120° with respect to each other. The rms value of neutral current is \hfill{[EE-2019]}
    \begin{multicols}{4}
        
    
    \begin{enumerate}
        \item 0 $A$
        \item $\frac{100}{\sqrt{3}}$ $A$
        \item 100 $A$
        \item 300 $A$
    \end{enumerate}
    \end{multicols}
    \item If $A = 2x i + 3y j + 4z k$ and $u = x^2 + y^2 + z^2$, then $\text{div}(uA)$ at (1, 1, 1) is \underline{\hspace{2cm}}. \\.\hfill{[EE-2019]}




\end{enumerate}



\end{document}


